\section*{Ringkasan}
\addcontentsline{toc}{section}{Ringkasan}

Sistem listrik Jawa-Bali menyerap 69,5\% penjualan listrik nasional \cite{D-PLN-STAT-2025} dan masih bertumpu pada 57 unit PLTU batu bara terhubung jaringan PLN dengan kapasitas 30,8 GW dan emisi sekitar 169 juta ton CO$_2$ per tahun \cite{D-GEM-GCPT}. Laporan Kelompok 17 memilih integrasi PLTS atap dan pensiun dini PLTU sebagai solusi transisi \cite{L-LAPORAN-K17}. PLTU dipensiunkan bertahap, kapasitasnya digantikan PLTS atap, beban investasi disebar melalui insentif dan kewajiban pemasangan, dan pekerja terdampak dilindungi.

Masalahnya bukan ketiadaan aturan. Kebijakan PLTS atap dan pensiun dini berjalan terpisah, sehingga PLTS atap dipandang PLN sebagai pengurang pendapatan, sedangkan pensiun dini dipandang sebagai beban kompensasi. Pembatalan pensiun dini Cirebon-1 pada Desember 2025 menunjukkan bahwa tanpa hitungan yang jelas tentang siapa menanggung apa, pensiun dini mudah ditunda \cite{B-BISNIS-CIREBON-BATAL-2025}.

Dokumen kajian ini merekomendasikan \textbf{program pasangan kuota PLTS atap dan pensiun dini PLTU IPP}. Setiap PLTU yang dipensiunkan dipasangkan dengan tambahan PLTS atap paling sedikit setara dayanya, didahului pengikatan dana luar, pembangkit pengganti, dan program bagi pekerja. Urutan pensiun mendahulukan PLTU milik IPP karena hanya pensiun PLTU IPP yang menghentikan pembayaran kapasitas oleh PLN. Seluruh penilaian dilakukan dengan platform terbuka PowerShift yang dikembangkan bersama kajian ini.

Simulasi roadmap Fase 1-3 di Jawa-Bali memensiunkan 3,0 GW PLTU IPP pada 2030-2035 dan menambah 3,19 GWp PLTS atap, atau 106\% dari daya yang dipensiunkan. Selama 2026-2040 paket ini menghindari 9.572 kematian dini dan 35,4 juta ton CO$_2$. Biaya bersih bagi negara Rp11,9 triliun, atau kurang dari Rp1 triliun per tahun, sementara pelanggan menghemat tagihan Rp38,0 triliun dan nilai kesehatan yang terselamatkan Rp111,3 triliun. Bila neraca negara harus nol, varian dengan PLTS atap lebih kecil menghasilkan surplus Rp1,8 triliun.

\begin{center}
\small
\begin{tabularx}{\textwidth}{*{4}{>{\centering\arraybackslash}X}}
\toprule
\textbf{3,0 GW} & \textbf{106\%} & \textbf{9.572 jiwa} & \textbf{Rp15,2 T} \\
PLTU IPP pensiun dini 2030-2035 & PLTS atap tambahan terhadap daya PLTU yang dipensiunkan & kematian dini dihindari 2026-2040 & dana luar, di bawah alokasi CIPP JETP sekitar Rp21,5 T \\
\bottomrule
\end{tabularx}
\end{center}
